
\clearpage

\section{Evaluating culture, inclusion and belonging}
\label{sec:culture}

\subsection{Provide information on how the department ensures that culture and practices support inclusion and belonging. This should include, but is not limited to, information on how the department actively considers gender equality, and EDI more broadly, in:} 
\instr{
\begin{itemize}
\item	organisation of meeting and events;
\item	images and text used in department spaces and on the department’s website; 
\item	student curricula, pedagogy, and assessment. 
\end{itemize}
}

\subsubsection{Organisation of meetings and events}
The school of CSIT organises a variety of school meetings, committees, social events and research seminars. %to support the diverse activities occurring at the school. 
Social events are normally organised for staff at Christmas and in summer. Staff are invited to register their time constraints  (w.r.t. family and caring responsibilities). 
%Other social events include a monthly coffee morning at 10:30-11:30, which continued on Microsoft Teams during the Covid-19 pandemic. 
%The coffee mornings continued weekly during the Covid-19 pandemic on Microsoft Teams to facilitate inclusion while working remotely.

Staff meetings take place on the last Friday of the month during the academic year at 15:00. The timing of the meetings is currently being discussed to ensure that the meeting ends within the core working hours (i.e. before 16:30); some topics might provoke debate among staff leading to meetings ending after 5pm. While it is important for everyone to have an opportunity to speak, the School realises that meeting exceeding core working hours are in effect negating this chance to staff with caring responsibilities.

The School has eight committees to support school activities and diversity. 
%Most academic and PMS staff are, or can become, a member of these committees (see previous discussion on academic workload on page \pageref{sub:workload}). 
%Membership is allocated by the Head of School alongside committee chairs, taking into consideration staff preferences (interests, career development, etc.). In 2021, the Head of School encouraged staff to provide preferences committee membership and roles.
Table~\ref{tab:gender:committees} shows 
%provides an overview of gender representation in school committees. There is 
a significant under-representation of female committee chairs, which is related to a lack of female academic staff in senior roles in CSIT.



\begin{tabbox}[11cm]{!th}
\caption{Gender representation in School committees}
\label{tab:gender:committees}
\sisetup{
    group-digits=all,
    group-minimum-digits=4,
    group-separator = {,},
    table-format=4.1,
    mode=text,
    reset-text-series = false, 
    text-series-to-math = true,
    reset-text-family = false, 
    text-family-to-math = true
}
\footnotesize   
\begin{tabular}{p{6cm} 
                !{\quad}
                S[table-format=2.0]
                !{\quad}
                S[table-format=2.0]
                !{\quad}
                S[table-format=2.0\%]
                !{\quad}
                S[table-format=2.0]
                }

\toprule 
Role & {Male} & {Female} & {Female \%} & {Total}\\
\midrule 
Committee Chairs & 8 & 0 & 0\% & 8\\
\hdashline
Committee Members & 16 & 7 & 30\% & 23 \\
\hdashline
All Committee (Chairs and Members) & 24 & 7 & 23\% &31\\
\bottomrule 
\end{tabular}


\end{tabbox}

The School hosts a seminar series to host both internal and external speakers. Table~\ref{tab:seminars} shows an under-representation of female speakers. Recently, the School started to make specific efforts to target female  speakers when issuing invitations.
All staff are invited to suggest speakers; when suggesting male speaker, staff are strongly requested to also suggest a female speaker.
%the more general gender imbalance in computer science internationally impacts the female representation among seminar speakers in the School. 
%The seminar series did not take place during 2021 due to the Covid-19 pandemic and was relaunched in January 2022. 

\begin{tabbox}[9cm]{!th}

\caption{Gender distribution of speakers at CSIT seminars}
\label{tab:seminars}
\sisetup{
    group-digits=all,
    group-minimum-digits=4,
    group-separator = {,},
    table-format=4.1,
    mode=text,
    reset-text-series = false, 
    text-series-to-math = true,
    reset-text-family = false, 
    text-family-to-math = true
}
\footnotesize   
\begin{tabular}{p{3cm} 
                S[table-format=2.0]
                !{\qquad}
                S[table-format=2.0]
                !{\qquad}
                S[table-format=2.0]
                !{\qquad}
                S[table-format=2.1\%]
                }

\toprule 
Year & {Total} & {Male} & {Female} & {Female \%}\\
\midrule 
2019 & 18 & 13 & 5 & 27.8\%\\
\hdashline
2020 & 11 & 9 & 2 & 18.1\%\\
\hdashline
2021 & \multicolumn{4}{c}{No seminars due to Covid-19} \\
\hdashline
2022 & 13 & 8 & 5 & 38.4\%\\
\midrule 
Total & 42 & 30 & 12 & 28.6\%\\
\bottomrule 
\end{tabular}


\end{tabbox}

%While academic and support staff are involved in the running of the school, 
Feedback from students and the research staff focus group revealed that research staff and students lamented a lack of events and opportunities where PhD students can engage with staff and the School. 
%The research centres hosted by the School (see Section~\ref{sec:overview}) have separate communication channels to their respective members, and events are frequently not communicated to other postgraduate students. 
Most of the events for networking, job opportunities or professional development are organised by the various research centres hosted by the School (see Section~\ref{sec:overview}), each of which has its own dedicated communication channels such as mailing lists; while these events are open to anyone, they are not clearly communicated to those students who are not members of the research centres. 
%Students not associated to one of the centres have very limited opportunities to interact. Very few events for networking, job opportunities or professional development events are organised explicitly by the school and communicated to postgraduate students. Most are organised by the different centres with them being open to students from other centres as well, but most often this is not communicated.
This in turn weakens these researchers' sense of belonging, limiting the opportunities for a shared work culture to develop among them. %and a work culture has little opportunity to develop between all student. 
PhD students are allocated a workspace in access-controlled laboratories primarily based on the area they work in, rather than the centre they belong to. While this allows opportunities for research discussions and collaboration, it limits interaction among students working in different areas. %But since access is controlled, and generally granted only to the lab the student belongs to, collaborating and discussing with students in other fields is more difficult.

While the University offers a Research Orientation, research students also lamented a lack of orientation offered by the School.  Little information is offered to newcomers regarding the different centres, the facilities, the School, or opportunities to get involved in teaching and lab assistance.  %demonstrating and other teaching-related topics. It was felt that this does not promote an environment of inclusivity. 
%that there is little introduction given to new students by the School. 
These issues hinder a culture of inclusivity within the school, and also affect research, as it partially negates a culture of multi-disciplinary collaboration. 
Thus we propose Action~\ref{action:acad:fora} that seeks to address this:
%For this reason we propose a related action aimed at creating events and spaces where such collaborations can occur.

\newcommand{\actacadfora}{%Action 2.13 - talk to Paolo about it
Create regular fora, meetings and social events for all academic and research staff (across research groups and centres) to make the school more inclusive for all staff.}
\begin{actionblock}{\ksspace{}Create Fora for Inclusivity}{acad:fora}
\actacadfora
\hfill{
    \hypersetup{hidelinks}
    \hyperlink{c9}{\gotoactionsymbol{}}
}
\end{actionblock}


%Not many policies exist for postgraduate students specifically. Most policies that concern students assume a student has responsibilities only towards courses/exams which take place in the academic year, not research. In most cases regulations/opportunities are coming from supervisor and funding agency.

\subsubsection{Images and text used in department spaces and on the department's website}
The School has its own dedicated website (\url{https://www.ucc.ie/en/compsci/}) to promote its programmes, research, outreach and all staff activities. There is a dedicated news section to promote school events, including PhD graduations and defences, undergraduate events, research, and outreach events, all of which include pictures of those involved.

The School invites students from all programmes to share their thoughts on our course offerings, which are advertised on the school website. Currently, the site includes 5 testimonials of our undergraduate programmes (3 female; 2 male).

\subsubsection{Student curricula, pedagogy, and assessment}
The School's four undergraduate programmes capture a diversity of interests  beyond computer science and technology: humanities, data and mathematics, psychology/human computer interaction. All programmes provide a range of pathways to cater for individual learning and interests, as shown in Table~\ref{tab:programmeoptions}.

\begin{tabbox}{!th}
\caption{Programme options}
\label{tab:programmeoptions}
\sisetup{
    group-digits=all,
    group-minimum-digits=4,
    group-separator = {,},
    table-format=4.1,
    mode=text,
    reset-text-series = false, 
    text-series-to-math = true,
    reset-text-family = false, 
    text-family-to-math = true
}
\footnotesize   
\begin{tabular}{p{5cm} 
                p{9cm}
                }

\toprule 
Programme & Options \\
\midrule 
BSc Computer Science & Computer Science, Maths and language electives\\
\hdashline
BSc Data Science and Analytics & Maths and Statistics electives\\
\hdashline
BA Psychology and Computing & Optional 9-month work placement\\
\hdashline
BA Digital Humanities and IT & Optional 9-month work placement; pursue an arts minor subject (e.g. languages, sociology, cultural studies)\\
\bottomrule 
\end{tabular}


\end{tabbox}

The school has two committees that are concerned with curricula, pedagogy, and assessment: \ac{APCD} committee and the \ac{TLSE} committee.
The APCD is responsible for reviewing changes to programmes and modules delivered within the school, as well as reviewing proposals for new academic programmes supported by the \ac{HOS}. The \ac{TLSE} is focused on student experience, including drafting policies and procedures to preserve and enhance the working environment, teaching and learning experience. 

The School's outreach events are geared toward supporting gender diversity, and include secondary school visits and participation in the annual iWish event for women in science.\footnote{\url{https://www.iwish.ie}} During secondary school visits, the School targets a diversity of prospective for male and female students, sending targeted emails to girls schools. 
%Female staff are also invited to present at Schools, if interested. 
Table~\ref{tab:schoolvisits} presents the number of school visits between 2018-2022. 
%In 2022, the Covid-19 pandemic had an impact on the schools visited. 
%Mixed schools saw the highest level of visits overall. 
While the School recognises the importance of visits to girls schools in particular, such school visits are dependent on the willingness of staff and the visited school to facilitate. 
%The table highlights there is an under-representation of girls schools visited.

\begin{tabbox}[10cm]{!th}

\caption{School visits to female-only, male-only, and mixed-gender schools}
\label{tab:schoolvisits}
\sisetup{
    group-digits=all,
    group-minimum-digits=4,
    group-separator = {,},
    table-format=4.1,
    mode=text,
    reset-text-series = false, 
    text-series-to-math = true,
    reset-text-family = false, 
    text-family-to-math = true
}
\footnotesize   
\begin{tabular}{p{3cm} 
                !{\qquad}
                S[table-format=2.0]
                !{\qquad}
                S[table-format=2.0]
                !{\qquad}
                S[table-format=2.0]
                !{\qquad}
                S[table-format=2.0]
                }

\toprule 
Year & {Female} & {Male} & {Mixed (M/F)} & {Total}\\
\midrule 
2018/19 & 1 & 4 & 3 & 8\\
\hdashline 
2019/20 & 2 & 1 & 2 & 5\\
\hdashline 
2020/21 (virtual) & 5 & 5 & 10 & 20\\
\hdashline 
2021/22 & 0 & 3 & 2 & 5\\
\midrule 
Total & 8 & 13 & 17 & 38 \\
\bottomrule 
\end{tabular}
\vspace{2mm}\\
{\sffamily
School visits during 2020/21 were virtual due to Covid-19; school visits did not resume in full during 2021/22.}\\



\end{tabbox}

All students are assigned an academic mentor, a member of academic staff who retains contact with the student and provides resources in cases of any issues with module material, pedagogy, and assessment.
For all programmes, students can apply for accommodations to support formal end-of-term exams (for example, through the university \ac{DSS}). Students in need of support for continuous assessment, are accommodated by the lecturer delivering the module. 
The school has no guidelines for supporting students who may need accommodations for continuous assessment at a school level.

While computer science is an inherently technical subject, the Athena SWAN process has increased awareness among staff on the necessity to introduce and discuss, where relevant, gender-related elements in their teaching, for example, by using gender-neutral examples.

\subsection{Comment and reflect on the department’s current understanding of, and capacity to identify and address, issues and opportunities relating to equality grounds in addition to gender, as well as capacity to identify and address intersectional inequalities for staff and students.}

While the School members come from varied ethnic backgrounds, the representation of different ethnicities varies by role, with no one identifying with a non-White background in managerial and professorial roles. 
%Staff hired in the last 10 years have had significantly reduced opportunities to progress in their career due to a lack of promotion calls, and therefore to obtain higher income, which in turn reduces opportunities for securing accommodation, with new staff frequently struggling to secure adequate housing due to the current housing crisis in Ireland. 
The lack of promotion calls and opportunities discussed earlier has economic consequences as well. A lack of promotions means that staff are not able to secure better wages, which in turn reduces opportunities to secure accommodation. The ongoing housing crisis in Ireland affects staff who are renting and who may not be in a position to buy a house; new staff frequently struggle to secure adequate accommodation.
%This has disproportionately affected female members of staff, who are, for the most part, relatively recent hires. 
The lack of job security in a number of positions, from postdoctoral researcher to Lecturer, also has socioeconomic implications that transcend and complement those of gender. Limited earning potential, combined with lack of adequate support and facilities for childcare, significantly impacts all staff hired in positions of Lecturer and above, but is specifically felt by international staff members with limited local support from extended family.

\subsection{Provide information on the department’s culture as it relates to gender equality and, where relevant, EDI more broadly, by presenting consultation findings by gender and staff category on the following areas:}
\instr{
\begin{itemize}
\item	values and traditions of the department;
\item	formal and informal structures and interactions that characterise the working and learning environment of the department, including leadership practices and behaviours;
\item	negative practices and behaviours and how these are managed by the department; 
\item	flexible working opportunities in the department;
\item	management of, and attitudes towards, family leave in the department. 
\end{itemize}
\vspace{2mm}
\noindent Where data suggests opportunity for improvement, comment and reflect. This should include reflection on any gaps between institution-level policy and practice in the department, including if the institution’s approach meets the requirements of department staff.
}

\subsubsection{Values and traditions of the department}
There is a strong collegiate culture within the School. 
All staff members are invited to social activities organised by the school, including monthly coffee meetings, the Christmas event, and the summer barbeque event. 

The School also facilitates various social events to get students and/or researchers together. For example, several online quiz nights were organised for postgraduate students during the pandemic. The School organises an annual reception for international students that is attended by the \ac{HOS}, Programme coordinators, and staff from the university's International Office. Furthermore, there are initiatives by individuals within the School; for example, one of the postdoctoral researchers organises weekly semi-social meeting to get researchers together, and one of the lecturers organised a well-attended PhD climbing wall event.

\subsubsection{Formal and informal structures and interactions characterising the department's working and learning environment}

All newly hired staff are assigned mentors to provide guidance on various matters to ensure that they are well-integrated.

To address gender imbalance among students, the School considers numerous practices to encourage female students to join the school. Examples of these practices include ensuring equal gender representation at the School's Open Day booth, online recruitment campaigns, and asking female staff to attend visits to girls schools.

Beyond gender, the \ac{HOS} recently initiated strategies to support minorities. For example, the School has recently instantiated an international student mentor role to advise international students on various academic and non-academic matters. 

The School considers numerous practices to ensure that everyone is integrated and heard. School meetings avoid non-core hours to ensure that everyone is able to attend. 
%Additionally, regular staff meetings are scheduled to ensure that staff have no clashes with other duties. 
These meetings are chaired by a senior member of staff who is not \ac{HOS}, who solicits agenda items from all staff. 
During the meetings, staff are offered the opportunity to speak their minds on all school matters. 


The School supports female staff and students for targeted grants. Similarly, the School communicates female-targeted students' scholarships (e.g., Google stipend) and encourages  students to apply for them. The School also encourages female students to participate or become a team member in various competitions such as the Irish programming Olympiad and the International Olympiad for Informatics.  
The School also advertises female-targeted events (e.g., IEEE Women in Engineering) and encourage students and staff to participate in these events.  


\subsubsection{Negative practices and behaviors and their management by the department}
\todo{no information?}

%\subsubsection{Leadership practices and behaviours \todo{to be merged into other sections?}}

%The school leaders adopt various practices to strengthen the sense of belonging and inclusion. 
%The following paragraphs highlight common actions and practices by school leaders for staff and student recruitment, supporting existing and students, and an enjoyable social life. 
%It is important to note that the school recognise gender imbalance among academic staff and students and continuously considers actions to close that gap. 
%The school was proactive in its engagement with the Athena SWAN process, and appointed a Champion to liaise with the university to get an understanding for Athena SWAN (2 years before a planned application) and fully committed to proceed with the award application.

%In the following, the leadership practices discussed in one-to-one interviews to (current and former) heads and managers are presented.

%\begin{itemize}

%\item	\textbf{Staff and Student Recruitment}: 
%The school leaders recognise gender imbalance among academic staff and have taken many steps to overcome. 
%HR was consulted to ensure that job advertisements include inviting statements for females. 
%Additionally, job openings are directly communicated to qualified candidates from various institutions. 
%These actions increased the number of female applicants for recently advertised positions leading to hiring a first female lecturer in 2017, and three additional female lecturers joined the staff in 2018 and 2019. 
%It is worth noting that some positions were offered to short-listed female candidates but they declined the offer due to the high teaching load. 
% The school considers evolving opportunities to overcome the gender imbalance among academic staff. 
% For example, the school tried twice to hire female professors in HCI and AI through the \ac{SALI} scheme but that was not successful. On hiring researchers, the directors of the research centres target maintaining gender balance. A similar goal is considered when recruiting Principal and Funded Investigators for research centres. 


%\item	\textbf{Existing Staff and Students}: 

%The \ac{HOS} follows a democratic process for load assignment. 
%Staff are surveyed for their preference (teaching and service tasks), and the assignment prioritises such preference, while ensuring meeting other objectives, such as balancing gender, type -of-job, and seniority level when forming committees. The school head and administration keep close contact with year heads to ensure that student voice is heard.  


%\item \textbf{Support Minorities.}	
%\end{itemize}


\subsubsection{Flexible Working}

The school currently supports flexible working for all staff. 
%Staff who take long-term leave (e.g., maternity or illness) are offered a reduced load upon their return. 
%\todo{Beyond the School remit, the current \ac{HOS} advises companies and IDA to adopt accommodating policies (e.g., flexible working and part-time) to encourage females to join and remain in work force.}

The majority of staff are allowed to work from home under the Pilot Blended Working Policy, however some staff have not been allowed due to the nature of their job.
Blended Working varies from academic to support staff due to the nature of the roles. The feedback from the PMS focus group as well as the EDI survey indicated a strong support for blended working to continue (75\% either agree or strongly agree to blended working) and participants expressed a need for clarity on this.
In the PMS focus group, staff expressed near-unanimous strong support for blended working, with this reflected in Action~\ref{action:pms:flexworking}. Staff acknowledged that the blended work policy may need tailoring to the role’s responsibilities but that every effort should be made to seek a solution that offers flexibility to staff. Many PMS staff members highlighted the need for the University to clarify future policy. 
%Others felt strongly that flexible start/finish times should be formalised, assuming the \textit{``job can still be done.''} 

\newcommand{\actpmsflexworking}{Ensure PMS line managers meaningful engage with PMS staff to ensure every opportunity is given to adopt flexible working practices in line with the function of their role. Denied requests should be afforded good reasoning. }

\begin{actionblock}{\ksspace{}Engage Meaningfully with PMS Staff on Flex Work}{pms:flexworking}
\actpmsflexworking
\hfill{
    \hypersetup{hidelinks}
    \hyperlink{d9}{\gotoactionsymbol{}}
}
\end{actionblock}

 

Working hours vary by role. The vast majority of staff said they felt comfortable discussing flexible working arrangements with their line manager. The School's reception is open daily 9:00 to 17:00 which limits some PMS staff to avail of flexible start and finish times. Most working arrangements seem to be informally arranged with line managers where applicable.  
%The ability to change one's schedule again depends on the role.
Academics can to some extent indicate a preference for certain time slots for scheduled laboratory sessions, though various other restrictions such as the availability of adequately sized rooms must also be considered. 
%For academic staff, changes to teaching schedules are facilitated where possible when asked but it should be noted that some restrictions can apply based on availability of adequately sized rooms/labs and on programmes taking certain modules. 

% Some other existing UCC policies which promote flexible working are as follows:
% \begin{itemize}[noitemsep]
% \item	Flexible Working Hours
% \item	Incentivised Shorter Working Year Scheme
% \item	Reduced Working Week
% \end{itemize}

%It is unclear from the Staff Survey or from the Focus Groups if anyone in the School is availing or has availed of the above. 

%Annual leave and when it can be taken should also be considered. 
Some restrictions to annual leave apply due to the nature of some roles; for example, academics cannot normally take annual leave during teaching periods. Similar restrictions apply to some PMS staff. 

Research students are given the chance for flexible working hours, if their supervisor approves. There is also the flexibility to work from home or abroad (such as the home country of international students) as long as the supervisor approves.

\subsubsection{Management and attitudes towards family Leave}

Family leave policies are stipulated and enacted centrally by the university, with the School  in charge of implementing the cover requirements for staff on leave. 
%While the School has limited flexibility in modifying the policies, certain criticalities emerged during 
During the Athena SWAN consultation process, one problematic issue was highlighted:
%In particular, 
UCC currently maintains a requirement for staff members to have a minimum of six months of service before they are eligible for paid maternity leave. This requirement puts UCC at odds with other Irish universities, where this requirement was either never in place, or was recently removed. For example, University College Dublin removed this requirement in a policy revision in February 2019. %where the service requirement to receive paid maternity leave was removed in a policy revision in February 2019.




\newcommand{\actcultmaternity}{Advocate a removal of the minimum service requirement (6 months) for staff to be eligible for paid maternity leave. Liaise with the University management to achieve its removal from the policy.}
\begin{actionblock}{\ksspace{}Advocate Removal Minimum Service}{cult:maternity}
\actcultmaternity
\hfill{
    \hypersetup{hidelinks}
    \hyperlink{e3}{\gotoactionsymbol{}}
}
\end{actionblock}

Participants  of the academic focus group perceived a supportive environment for colleagues taking maternity and paternity leave. A participant mentioned being aware of a colleague who was offered a teaching load reduction on returning from maternity leave, which they felt was a positive way of supporting her \textit{``to kind of get back into it after being off.''}   Participants commended both the current and previous Heads of School for their support for leave-takers. %, including in the support and arrangement of backfill cover. 

%While backfill and cover arrangements have been provided, it 
With respect to some forms of family-related leave, including paternity leave and parental leave, the provision of backfill is not formally required, \textit{``so it's up to the Head of School to try and find a way to cover your classes [and] reduce [your] teaching load for that semester.''}   
Without a corresponding institutional commitment to support the arrangement of backfill/cover in these circumstances, Heads of School carry a responsibility that some participants felt may be unreasonable.  It also leaves leave-takers dependent on the goodwill and understanding of their particular \ac{HOS} and puts an increased demand on the goodwill of colleagues, who might be asked to help cover work during a period of leave.  
Leave-takers can feel a sense of responsibility and guilt for burdening colleagues with increased workload; this in turn might deter people from taking statutory leave:    

% \begin{myquote}
% “I was lucky because they had the cover for me and [there was help from] other members of staff. If [other people] don't have that, who knows?” 
% \end{myquote}

\begin{myquote}
“You feel bad about [taking leave] because you're basically asking [colleagues] to [cover] [your] classes ... It's very awkward, and that has nothing to do with the willingness of the Head of School. As I said, [they’re] very supportive, but it's all managed informally. They have to try and help you without getting any [support themselves], pretty much.”  
\end{myquote}

\newcommand{\actcultbackfill}{Formalising cover and backfill arrangements for family leave other than maternity leave (including parent’s leave, parental leave, paternity leave). Remove the expectation that the person taking leave should arrange for cover and backfill.}
\begin{actionblock}{\ksspace{}Formalise Cover and Backfill}{cult:backfill}
\actcultbackfill
\hfill{
    \hypersetup{hidelinks}
    \hyperlink{e4}{\gotoactionsymbol{}}
}
\end{actionblock}



Regarding research students, current policy states that a pregnant student should notify the supervisor and independently identify options she has given her funding. There is no specific UCC support. %, it is all depending on the student's funding body and the supervisor. 
The student can take a leave of absence, but should first make sure of the consequences of this in her exact situation. There is no suitable university housing to support students with children.  

The cr\`{e}che services offered by the university are limited, and while both staff and students are eligible, students are given priority. Due to the limited capacity,  only children of students have been admitted for the past two years. Research students are eligible under the same rules as staff, however the cr\`{e}che is limited to the academic year, whereas research students work fulltime throughout the year.

International PhD students who move to Ireland obtain a student visa, with which bringing over their partner/family can be difficult as the visa they obtain will not allow the partner to work. This puts such families in a financially difficult position given the low level of a PhD stipend. 





\newcommand{\actcultdiversity}{Increase outreach to students with diverse backgrounds. In particular, under-represented communities and ethnic backgrounds (e.g. travellers, Asian) will be targeted. Seek mechanisms to support Refugees and Asylum seekers.}
\begin{actionblock}{\ksspace{}Increase Diversity}{cult:diversity}
\actcultdiversity
\hfill{
    \hypersetup{hidelinks}
    \hyperlink{e1}{\gotoactionsymbol{}}
}
\end{actionblock}


\newcommand{\actcultmeetinghours}{Review hours of meetings in line with university policies on core working hours (10:00 to 16:30) including School Council. Review opening hours of the School Admin Office Counter. Facilitation of blended working within the School, and continuation of the programme at University level.}
\begin{actionblock}{\ksspace{}Review Meeting Hours}{cult:meetinghours}
\actcultmeetinghours
\hfill{
    \hypersetup{hidelinks}
    \hyperlink{e2}{\gotoactionsymbol{}}
}
\end{actionblock}





\newcommand{\actcultmatfunding}{Advocating with research funders and the University to achieve a single formal scheme/arrangement for the payment of and provision of funding for maternity and other family leave to research students and staff funded on research projects/grants.}
\begin{actionblock}{\ksspace{}Advocate Funding Maternity Leave}{cult:matfunding}
\actcultmatfunding
\hfill{
    \hypersetup{hidelinks}
    \hyperlink{e5}{\gotoactionsymbol{}}
}
\end{actionblock}

